\documentclass[conference, a4paper]{IEEEtran}
\IEEEoverridecommandlockouts

\usepackage{cite}
\usepackage{amsmath,amssymb,amsfonts}
\usepackage{graphicx}
\usepackage[font=footnotesize]{caption}
\usepackage{textcomp}
\usepackage{xcolor}
\usepackage{url}
\usepackage{booktabs}
\usepackage{array}
\usepackage{geometry}
\usepackage{titlesec}

\geometry{
  a4paper,
  top=2cm,
  bottom=2.5cm,
  left=1.6cm,
  right=1.6cm,
}

\setlength{\parskip}{6pt}
\setlength{\floatsep}{0pt}
\setlength{\textfloatsep}{6pt}
\setlength{\intextsep}{6pt}
\titlespacing*{\section}{0pt}{8pt}{4pt}
\titlespacing*{\subsection}{0pt}{6pt}{3pt}
\titlespacing*{\subsubsection}{0pt}{4pt}{2pt}

\begin{document}

\title{Hierarchical Deep Learning for Carbonate Grain Classification in Thin-Section Images}

\makeatletter
\def\@IEEEauthorblockNfont{\fontsize{11pt}{13pt}\selectfont}
\def\@IEEEauthorblockAfont{\fontsize{10pt}{12pt}\selectfont}
\makeatother

\author{
\IEEEauthorblockN{Musawer Muradi}
\IEEEauthorblockA{\textit{Department of Computer Engineering}\\
\textit{Kocaeli University}\\
Kocaeli, Turkiye\\
musavir.muradi@gmail.com}
\and
\IEEEauthorblockN{Arnaud Gallois}
\IEEEauthorblockA{\textit{Department of Geological Engineering}\\
\textit{Kocaeli University}\\
Kocaeli, Turkiye\\
arnaud.gallois@kocaeli.edu.tr}
\and
\IEEEauthorblockN{Alev Mutlu}
\IEEEauthorblockA{\textit{Department of Computer Engineering}\\
\textit{Kocaeli University}\\
Kocaeli, Turkiye\\
alev.mutlu@kocaeli.edu.tr}
}

\maketitle

\begin{abstract}
Carbonate grain classification from thin-section images is difficult because the visual differences between classes are often small and the class distribution is strongly imbalanced. In the dataset used here, 2,642 annotated grains from 18 plain polarised light micrographs are grouped into four classes: peloid, ooid, broken ooid, and intraclast. Peloids account for 87.1\% of the grains, while broken ooids account for only 1.1\%. This paper presents a hierarchy-aware ResNet-18 model that follows the geological decision structure of the task: peloid versus non-peloid, ooid-like grain versus intraclast, and intact ooid versus broken ooid. Each decision head is trained with stage-specific Focal Loss, targeted oversampling is applied to broken ooids, and a staged training schedule is used to reduce majority-class dominance. On a held-out test set of 529 grains, the full model achieves 83.5\% Balanced Accuracy and 92.1\% Overall Accuracy. Broken-ooid recall increases from 16.7\% in the hierarchical model without oversampling to 83.3\% in the full configuration.
\end{abstract}

\renewcommand{\IEEEkeywordsname}{Keywords}
\begin{IEEEkeywords}
\itshape carbonate microfacies, deep learning, class imbalance, hierarchical classification, Focal Loss
\end{IEEEkeywords}

\section{Introduction}
Carbonate thin-section interpretation is a routine but demanding part of sedimentological analysis. Grain types such as peloids, ooids, intraclasts, and broken ooids are identified from subtle morphological and textural features. The task is not always reproducible, even for trained observers, because diagenetic alteration can hide the original grain texture. In the material studied here, micritisation is especially important: original grain fabrics may be replaced by fine calcite crystals, making distinct grains appear visually similar or transforming them into peloids \cite{bathurst1971,flugel2010}.

These visual ambiguities make carbonate grain classification a difficult image-recognition problem. Deep learning has shown promise in petrographic and geological image analysis \cite{koeshidayatullah2020,pireslima2020,liu2020}, but carbonate microfacies differ from many macroscopic rock-classification tasks. The target classes do not differ by large, obvious visual patterns. Instead, the useful signal may be a faint cortex, a fractured margin, or a remnant internal fabric.

A second difficulty is the class distribution. In the dataset used here, peloids account for 87.1\% of annotated grains, whereas broken ooids represent only 1.1\%. A model trained directly on this distribution can obtain high overall accuracy while failing on the minority classes \cite{johnson2019imbalance}. This is a poor outcome for geological use, because rare grains can still carry useful information and should not be treated as noise.

We therefore frame carbonate grain recognition as a hierarchy of binary decisions rather than a flat four-class problem. The proposed model uses a shared ResNet-18 backbone and three decision heads that follow the classification logic used by sedimentologists. The main contribution is not only the architecture, but the combination of hierarchy, stage-specific loss weighting, targeted oversampling, staged training, and inference-time augmentation under a severe imbalance ratio.

\section{Dataset}
The dataset contains 2,642 grain annotations from 18 plain polarised light micrographs of carbonate thin sections from the Mupe Member of the Purbeck Limestone Group in Dorset, southern UK. The samples are described in previous geological studies of the same succession \cite{gallois2016,gallois2018}. Grain boundaries were annotated with LabelMe \cite{labelme} by an expert sedimentologist, following standard carbonate petrographic criteria \cite{scholle2003carbonate,flugel2010}.

Four classes are considered: peloid, ooid, broken ooid, and intraclast. Other sparse grain types, including ostracod fragments, bivalve fragments, and quartz grains, were excluded from the classification task. Each annotated grain was cropped into a 96$\times$96 RGB patch centred on the annotation centroid. A binary mask from the polygon annotation was applied so that pixels outside the grain boundary were set to zero. This makes the classifier focus on the grain interior and boundary rather than the surrounding cement or matrix.

\begin{table}[t]
\caption{Class distribution in the annotated dataset.}
\label{tab:class_distribution}
\centering
\begin{tabular}{lrr}
\toprule
Class & Count & Percentage \\
\midrule
Peloid & 2300 & 87.1\% \\
Ooid & 210 & 7.9\% \\
Intraclast & 103 & 3.9\% \\
Broken ooid & 29 & 1.1\% \\
\midrule
Total & 2642 & 100.0\% \\
\bottomrule
\end{tabular}
\end{table}

The data were split into training, validation, and test sets using stratified random sampling with a 60/20/20 ratio and a fixed seed. The test set contains 529 grains: 460 peloids, 42 ooids, 6 broken ooids, and 21 intraclasts. It was held out during model development and threshold selection.

\begin{figure}[t]
\centering
\includegraphics[width=\linewidth]{figs/fig2_sample_patches.pdf}
\caption{Representative masked 96$\times$96 grain patches for the four classes: (A) peloid, (B) ooid, (C) broken ooid, and (D) intraclast. White scale bars indicate 50~$\mu$m. The visual overlap between heavily micritised grains is the central difficulty of the task.}
\label{fig:patches}
\end{figure}

\section{Method}
\subsection{Hierarchical Model}
A flat classifier maps each patch directly to one of four classes. Under the present imbalance, this encourages the network to favour the majority class. We instead decompose the problem into three binary decisions:
\begin{enumerate}
\item Stage 1: peloid versus non-peloid;
\item Stage 2: ooid-like grain versus intraclast;
\item Stage 3: intact ooid versus broken ooid.
\end{enumerate}

The backbone is ResNet-18 initialised with ImageNet weights \cite{he2016resnet,russakovsky2015imagenet}. The final classification layer is removed, producing a 512-dimensional feature vector for each masked grain patch. This vector is passed to three independent multi-layer perceptron heads. Each head contains a 512-to-128 linear layer, ReLU activation, dropout, and a 128-to-1 output layer followed by a sigmoid. In the implementation, Stage 2 uses ooid-like as the positive class and Stage 3 uses intact ooid as the positive class. The Stage 3 head is trained only on ooid-like grains, so it does not receive gradient updates from the large number of peloid samples.

\begin{figure}[t]
\centering
\includegraphics[width=\linewidth]{figs/architecture.png}
\caption{Hierarchical ResNet-18 model. A shared backbone feeds three binary heads that follow the carbonate grain classification sequence.}
\label{fig:architecture}
\end{figure}

\subsection{Loss, Sampling, and Training}
Each head is trained with Focal Loss \cite{lin2017focal}:
\begin{equation}
\mathcal{L}_{FL}(p_t)=-\alpha_t(1-p_t)^\gamma\log(p_t),
\label{eq:focal_loss}
\end{equation}
where $p_t$ is the predicted probability for the correct class, $\gamma$ controls the down-weighting of easy examples, and $\alpha_t$ provides class balancing. The balancing weight is set separately for each decision head: $\alpha=0.25$ for Stage 1, $\alpha=0.50$ for Stage 2, and $\alpha=0.75$ for Stage 3.

Loss reweighting alone does not solve the scarcity of broken ooids, because there are fewer than 20 broken-ooid examples in the training split. We therefore oversample broken ooids at 15 times their natural frequency. These are not synthetic grains: the same real grain patches are repeated with standard training augmentations, including random rotations, horizontal and vertical flips, and modest brightness and contrast changes.

Training follows four phases. First, the backbone is frozen and the three heads are initialised. Second, the backbone is unfrozen and tuned for Stage 1. Third, Stage 1 is frozen while Stages 2 and 3 are trained on oversampled batches. Finally, all parameters are fine-tuned jointly at a low learning rate. Training used AdamW with weight decay $10^{-4}$, batch size 32, dropout 0.3 in each head, and a fixed seed of 42. The four phases used 10, 15, 15, and 10 epochs with learning rates $10^{-3}$, $10^{-4}$, $10^{-4}$, and $10^{-5}$, respectively. Final routing used fixed thresholds of 0.5 at all three heads. At inference time, predictions are averaged over six label-preserving geometric transformations \cite{shanmugam2021tta}: the original image, horizontal flip, vertical flip, and rotations by 90, 180, and 270 degrees.

\section{Results and Discussion}
Balanced Accuracy \cite{brodersen2010balanced} is used as the primary metric because Overall Accuracy is misleading under this class distribution. A model that predicts peloid for every grain would already reach 87.1\% Overall Accuracy while being useless for the minority classes.

\begin{table*}[t]
\caption{Per-class recall on the 529-grain held-out test set for ablations of the proposed hierarchical pipeline. BA is Balanced Accuracy and OA is Overall Accuracy.}
\label{tab:ablation}
\centering
\resizebox{\textwidth}{!}{%
\begin{tabular}{lcccccc}
\toprule
Model & Peloid & Ooid & Broken ooid & Intraclast & BA & OA \\
\midrule
Hierarchical ResNet-18, no oversampling & 442 (96.1\%) & 33 (78.6\%) & 1 (16.7\%) & 13 (61.9\%) & 63.3\% & 92.4\% \\
+ 15$\times$ oversampling & 454 (98.7\%) & 35 (83.3\%) & 4 (66.7\%) & 13 (61.9\%) & 77.7\% & 95.7\% \\
+ 15$\times$ oversampling + TTA & 454 (98.7\%) & 37 (88.1\%) & 4 (66.7\%) & 14 (66.7\%) & 80.0\% & 96.2\% \\
+ Staged training + oversampling + TTA & 431 (93.7\%) & 36 (85.7\%) & 5 (83.3\%) & 15 (71.4\%) & 83.5\% & 92.1\% \\
\bottomrule
\end{tabular}
}
\end{table*}

The main effect comes from oversampling. Without it, the hierarchical model recovers only one of the six broken ooids in the test set. Adding 15$\times$ oversampling raises broken-ooid recall to 66.7\% and Balanced Accuracy from 63.3\% to 77.7\%. Inference-time augmentation gives a smaller but consistent gain, mainly by stabilising decisions for ooids and intraclasts. The full staged-training model reaches 83.5\% Balanced Accuracy and recovers five of the six broken ooids.

For context, a flat ResNet-18 trained with standard cross-entropy reached 86.9\% Balanced Accuracy and 95.5\% Overall Accuracy on the same test split. This flat model is reported as an external reference point, while Table~\ref{tab:ablation} focuses on ablations of the proposed hierarchical pipeline. Its higher score is driven mainly by correctly labelling all six broken ooid test examples, so it should be read cautiously: one changed prediction would move broken-ooid recall by 16.7 percentage points. The hierarchical model gives stronger intraclast recall, 71.4\% compared with 61.9\%, and its design is more aligned with the geological structure of the task.

The main remaining error is confusion between peloids and intraclasts. This is not only a model failure; it reflects a real petrographic ambiguity. When micritisation is advanced, both classes may appear as rounded grains with largely structureless micritic interiors. The model therefore struggles at the same boundary that can also be uncertain for human interpretation.

\section{Conclusion}
This study shows that, within a hierarchy-aware framework, targeted oversampling and staged training substantially improve minority-class recovery in carbonate thin-section image classification. The proposed ResNet-18 model follows the sedimentological decision sequence and combines stage-specific Focal Loss, targeted oversampling, staged training, and inference-time augmentation. On a held-out test set, the full model achieves 83.5\% Balanced Accuracy and 92.1\% Overall Accuracy, while broken-ooid recall increases from 16.7\% without oversampling to 83.3\%.

The results are still limited by dataset size and scope. Only six broken ooids are present in the test set, and all images come from a single carbonate succession. The next step is evaluation on independent thin-section datasets from other formations and imaging conditions. Even so, the results suggest a practical path for semi-automated carbonate petrography: use geological structure to guide model design, and treat rare grain types as central to the learning problem rather than as statistical outliers.

\section*{Acknowledgment}
The authors thank the contributors of the original petrographic material and annotations used in this study.

\bibliographystyle{IEEEtran}
\bibliography{references}

\end{document}
